% ---- score2 standard preamble (from: python3 score2_tools.py preamble) ----
\documentclass[10pt]{extarticle}
\usepackage[a4paper,margin=1.5cm]{geometry}
\usepackage[T1]{fontenc}
\usepackage{amsmath,amssymb}
\usepackage{graphicx}
\usepackage[export]{adjustbox}
\usepackage{multicol}
\usepackage{enumitem}
\usepackage{booktabs,array,tabularx}
\setlength{\columnsep}{0.7cm}
\setlength{\parindent}{0pt}
\setlength{\parskip}{3pt}
\pagestyle{empty}
\makeatletter
% figure registry (figdefs.tex is generated by the extract step)
\newcommand{\qfigdef}[4]{\@namedef{qf@file@#1}{#2}\@namedef{qf@w@#1}{#3}\@namedef{qf@wide@#1}{#4}}
\newcommand{\qf@check}[1]{\@ifundefined{qf@file@#1}{\PackageError{score2}{Unknown figure id #1}{}}{}}
\newcommand{\qf@img}[4][]{\includegraphics[width=#2\dimexpr\@nameuse{qf@w@#4}pt\relax,max width=#3,max height=0.45\textheight,#1]{\@nameuse{qf@file@#4}}}
% Full-width questions: finalize wraps every question that contains a wide figure in \qfullbegin...\qfullend,
% so the WHOLE question (stem, figures, options) is set across the page in reading order, then the two columns
% resume. The question's head - its number, stem and first figure (\qfullsplit marks the end of the head) -
% is kept on one page, so a number is never left at a page foot with its figure overleaf; the rest may break
% across pages normally, so no big gaps are left. Without a \qfullsplit, the whole question is kept together
% if it is under 45% of a page. Never write these yourself.
\newif\ifqf@full
\newif\ifqf@boxed
\newsavebox{\qf@box}
\newcommand{\qf@open}{\global\qf@fulltrue\par\global\qf@boxedtrue
  \setbox\qf@box\vbox\bgroup\hsize\textwidth\linewidth\textwidth}
\newcommand{\qf@emit}[1]{\par\egroup\global\qf@boxedfalse
  \ifdim\dimexpr\ht\qf@box+\dp\qf@box\relax<#1\textheight\noindent\usebox{\qf@box}\par\else\unvbox\qf@box\fi}
\newcommand{\qfullsplit}{\ifqf@boxed\qf@emit{0.85}\fi}
\newcommand{\qf@close}{\ifqf@boxed\qf@emit{0.45}\fi\global\qf@fullfalse}
\newcommand{\qfullbegin}{\end{multicols}\qf@open}
\newcommand{\qfullend}{\qf@close\begin{multicols}{2}}
\newcommand{\qfullstart}{\qf@open}
\newcommand{\qfulldone}{\qf@close}
% \qfig[scale]{id}: block figure, centred, directly where it is written (size set by the tool).
% Fallback if a wide figure is not inside a full-width question: break out of the columns in place - no floats.
\newcommand{\qfig}[2][1]{\qf@check{#2}%
  \ifqf@full\qf@blk{#1}{#2}\else
  \ifnum\@nameuse{qf@wide@#2}=1 \expandafter\@firstoftwo\else\expandafter\@secondoftwo\fi
  {\end{multicols}\begin{center}\qf@img{#1}{\textwidth}{#2}\end{center}\begin{multicols}{2}}%
  {\qf@blk{#1}{#2}}\fi}
\newcommand{\qf@blk}[2]{\par\smallskip{\centering\qf@img{#1}{\linewidth}{#2}\par}\smallskip}
% \qfigi[scale]{id}: inline figure (inside option lists / tables / minipages); never breaks columns.
\newcommand{\qfigi}[2][1]{\qf@check{#2}\qf@img[valign=c]{#1}{\linewidth}{#2}}
\makeatother
% structure macros - lint reads these to check completeness and figure placement
\newcommand{\papersec}[1]{\par\bigskip{\centering\Large\bfseries PAPER-#1\par}\medskip}
\newcommand{\partsec}[2]{\par\medskip\noindent{\large\bfseries PART-#1: #2}\par\smallskip}
\newcommand{\sect}[1]{\par\smallskip\noindent\textbf{#1}\par}
\newcommand{\qn}[1]{\par\smallskip\noindent\textbf{#1.}~}
\newcommand{\paraq}[2]{\par\smallskip\noindent\textbf{Paragraph for Questions #1 and #2}\par}
\newenvironment{qoptions}{\begin{itemize}[leftmargin=1.8em,itemsep=1pt,topsep=2pt,label={}]}{\end{itemize}}
\newcommand{\opt}[2]{\item[(#1)] #2}
% figure registry (generated by score2_tools.py finalize)
\qfigdef{p1_phy_s1_q01_1}{Score-I - Test 24 (11 Dec 2022)-img/p1_phy_s1_q01_1.png}{94.6}{0}
\qfigdef{p1_phy_s1_q03_1}{Score-I - Test 24 (11 Dec 2022)-img/p1_phy_s1_q03_1.png}{113.7}{0}
\qfigdef{p1_phy_s1_q04_1}{Score-I - Test 24 (11 Dec 2022)-img/p1_phy_s1_q04_1.png}{128.3}{0}
\qfigdef{p1_phy_s1_q05_1}{Score-I - Test 24 (11 Dec 2022)-img/p1_phy_s1_q05_1.png}{74.0}{0}
\qfigdef{p1_phy_s1_q06_1}{Score-I - Test 24 (11 Dec 2022)-img/p1_phy_s1_q06_1.png}{116.3}{0}
\qfigdef{p1_phy_s1_q07_1}{Score-I - Test 24 (11 Dec 2022)-img/p1_phy_s1_q07_1.png}{58.2}{0}
\qfigdef{p1_phy_s1_q09_1}{Score-I - Test 24 (11 Dec 2022)-img/p1_phy_s1_q09_1.png}{92.0}{0}
\qfigdef{p1_phy_s1_q10_1}{Score-I - Test 24 (11 Dec 2022)-img/p1_phy_s1_q10_1.png}{105.5}{0}
\qfigdef{p1_phy_s2_q02_1}{Score-I - Test 24 (11 Dec 2022)-img/p1_phy_s2_q02_1.png}{158.0}{0}
\qfigdef{p1_phy_s2_q03_1}{Score-I - Test 24 (11 Dec 2022)-img/p1_phy_s2_q03_1.png}{102.3}{0}
\qfigdef{p1_phy_s2_q04_1}{Score-I - Test 24 (11 Dec 2022)-img/p1_phy_s2_q04_1.png}{174.9}{0}
\qfigdef{p1_phy_s2_q07_1}{Score-I - Test 24 (11 Dec 2022)-img/p1_phy_s2_q07_1.png}{141.3}{0}
\qfigdef{p1_chem_s1_q01_1}{Score-I - Test 24 (11 Dec 2022)-img/p1_chem_s1_q01_1.png}{108.9}{0}
\qfigdef{p1_chem_s1_q01_A}{Score-I - Test 24 (11 Dec 2022)-img/p1_chem_s1_q01_A.png}{79.1}{0}
\qfigdef{p1_chem_s1_q01_B}{Score-I - Test 24 (11 Dec 2022)-img/p1_chem_s1_q01_B.png}{77.9}{0}
\qfigdef{p1_chem_s1_q01_C}{Score-I - Test 24 (11 Dec 2022)-img/p1_chem_s1_q01_C.png}{78.2}{0}
\qfigdef{p1_chem_s1_q01_D}{Score-I - Test 24 (11 Dec 2022)-img/p1_chem_s1_q01_D.png}{79.9}{0}
\qfigdef{p1_chem_s3_q01_1}{Score-I - Test 24 (11 Dec 2022)-img/p1_chem_s3_q01_1.png}{169.5}{0}
% ---- body skeleton ----
% \begin{document}
% \begin{center}{\Large\bfseries JEE (Main+Advanced) Score-II --- Test NN (DD Mon YYYY)}\end{center}
% \begin{multicols}{2}
%   \papersec{1}   (only for two-paper tests)
%   \partsec{1}{PHYSICS} \sect{Section-I (i): Single Correct} \qn{1} ... \qfig{p1_phy_q01_1} ...
% \end{multicols}
% \end{document}

\begin{document}
\begin{center}{\Large\bfseries JEE (Main+Advanced) Score-I --- Test 24 (11 Dec 2022)}\end{center}
\raggedcolumns
\setlength{\tabcolsep}{3pt}
\begin{multicols}{2}
\papersec{1}
\partsec{1}{PHYSICS}
\sect{SECTION-I(i)}
\qn{1} In the circuit shown in the figure, the value of $V_B+V_D$ :-
\qfig{p1_phy_s1_q01_1}
\begin{qoptions}
\opt{A}{$-3$V}
\opt{B}{7V}
\opt{C}{$-7$V}
\opt{D}{3V}
\end{qoptions}

\qn{2} Choose correct statement(s) :

(1) The number of significant figures in $0.31\times10^{3}$ m is two.

(2) The length and breadth of a rectangular sheet are 4.234 m and 1.00 m then the perimeter of rectangle with correct number of significant figures is 10.46 m.

(3) When 97.52 m is divided by 2.54 s, the result with correct number of significant figures is 38.393.

(4) When 0.2 J is subtracted from 7.26 J, then result with correct number of significant figure is 7.1 J.
\begin{qoptions}
\opt{A}{Only 1 and 2 are correct.}
\opt{B}{Only 1, 2 and 3 are correct.}
\opt{C}{Only 2 and 3 are correct.}
\opt{D}{Only 1 and 4 are correct.}
\end{qoptions}

\qn{3} In the circuit shown, the voltmeter reads 30 V when connected across the $400\ \Omega$ resistance. The same voltmeter when connected across $300\ \Omega$ will measure :
\qfig{p1_phy_s1_q03_1}
\begin{qoptions}
\opt{A}{20 V}
\opt{B}{22.5 V}
\opt{C}{25 V}
\opt{D}{40 V}
\end{qoptions}

\qn{4} The diagram shows two galvanometers $G_1$ and $G_2$. When current $I=1$A both $G_1$ and $G_2$ shows full scale deflection. It is given that $G_1$ shows full scale deflection for 5 mA and $G_2$ shows full scale deflection for 2 mA. The values of $r_1$ and $r_2$ are ($G_1$ and $G_2$ are of negligible resistances)
\qfig{p1_phy_s1_q04_1}
\begin{qoptions}
\opt{A}{$r_1=\dfrac{2R}{9}$ and $r_2=\dfrac{R}{995}$}
\opt{B}{$r_1=\dfrac{2R}{3}$ and $r_2=\dfrac{7R}{995}$}
\opt{C}{$r_1=\dfrac{3R}{2}$ and $r_2=\dfrac{7R}{995}$}
\opt{D}{$r_1=\dfrac{4R}{9}$ and $r_2=\dfrac{4R}{90}$}
\end{qoptions}

\sect{SECTION-I(ii)}
\qn{5} A block of mass `$m$' is attached by means of a spring to the bottom of a tank of water as shown in figure. At equilibrium, the spring is under compression. If the tank is now allowed to fall freely, then choose the correct alternative(s) :-
\qfig{p1_phy_s1_q05_1}
\begin{qoptions}
\opt{A}{the block comes to rest when the spring is at its natural length.}
\opt{B}{the spring compresses more than its equilibrium compression.}
\opt{C}{the buoyant force becomes zero.}
\opt{D}{there will be some elongation is the spring.}
\end{qoptions}

\qn{6} A piston is pressed with a force $F$ on a hydraulic cylinder containing water ($\rho=10^{3}$ kg/m$^3$). The cross-sectional area of the cylinder is $A=100$ cm$^2$. The water is forced into pipe with a cross-sectional area of 1 cm$^2$ that rises a height $h=50$ cm (Atmospheric pressure $P_0=10^{5}$ N/m$^2$)
\qfig{p1_phy_s1_q06_1}
\begin{qoptions}
\opt{A}{Force $F$ required to make water eject with a speed of 10 m/s is 550 N.}
\opt{B}{Force $F$ required to make water eject with a speed of 10 m/s at the top is 5500 N.}
\opt{C}{The pressure inside the cylinder at the bottom when water ejects at 10 m/s is $1.55\times10^{5}$ Pa.}
\opt{D}{none of the above}
\end{qoptions}

\qn{7} An isolated parallel plate capacitor have unequal charges $+4Q$ and $-2Q$ on its plates. If capacitance of the capacitor is $C$ and $P$ is a point outside of the capacitor. Then :-
\qfig{p1_phy_s1_q07_1}
\begin{qoptions}
\opt{A}{a point charge at point $P$ will not experience any electric force due to capacitor plates.}
\opt{B}{the potential difference between the plates will be $\dfrac{3Q}{C}$.}
\opt{C}{the energy stored in the electric field in the region between the plates is $\dfrac{9Q^2}{2C}$.}
\opt{D}{the force on one plate due to the other plate is $\dfrac{2Q^2}{Cd}$.}
\end{qoptions}

\qn{8} A cylindrical vessel $C_1$ of area of cross-section $A$ and height 1 m is filled fully up to a height 0.4 m by water (density 1000 kg-m$^{-3}$) and remaining by liquid of density 800 kg-m$^{-3}$. The cylinder has small orifice at the bottom, which is closed initially. The cylinder $C_1$ is then immersed in another bigger cylindrical vessel $C_2$ of area of crosssection $6A$, containing liquid of density 800 kg-m$^{-3}$ up to the height of 0.5 m. If the orifice of $C_1$ is opened, then,
\begin{qoptions}
\opt{A}{The speed with which the water comes out from the orifice (in ms$^{-1}$) will be 4}
\opt{B}{After certain interval of time, water from $C_1$ stops flowing out, thereby establishing the steady state of the system. The length of water column in cylinder $C_1$ is 6.67 cm}
\opt{C}{After certain interval of time, water from $C_1$ stops flowing out, thereby establishing the steady state of the system. The length of water column in cylinder $C_1$ is 9.5 cm}
\opt{D}{In steady state, the pressure exerted on the bottom of cylinder $C_2$ will be about $1.09\times10^{5}$ Nm$^{-2}$ [Atmospheric pressure, $p_0=1\times10^{5}$ Nm$^{-2}$ ]}
\end{qoptions}

\qn{9} In the circuit given below, if $R_1=240\ \Omega$ and $R_2=120\ \Omega$ respectively, L is a lamp of 300 W then
\qfig{p1_phy_s1_q09_1}
\begin{qoptions}
\opt{A}{Power dissipated in $R_1$ is 60 W}
\opt{B}{Power dissipated in $R_2$ is 120 W}
\opt{C}{Effective resistance of the circuit is $50\ \Omega$}
\opt{D}{The ammeter A reads 6 A}
\end{qoptions}

\qn{10} A battery of emf $E$ and internal resistance $r$ is connected across a resistance $R$. Resistance $R$ can be adjusted to any value greater than or equal to zero. A graph is plotted between the current passing through the resistance ($i$) and potential difference across the terminals of the battery ($V$). Select the correct alternative (s).
\qfig{p1_phy_s1_q10_1}
\begin{qoptions}
\opt{A}{Internal resistance of the battery is $4\ \Omega$.}
\opt{B}{emf of the battery is 16 V.}
\opt{C}{Maximum current which can be taken from the battery is 2A.}
\opt{D}{Maximum thermal power dissipated across $R$ is 16 J.}
\end{qoptions}

\sect{SECTION-III}
\qn{1} Find the attraction force (in Newton) between two parallel glass plates separated by a distance $h=0.10$ mm, after a water drop of mass $m=70$ g was introduced between them. The wetting is assumed to be complete. ($g=10$ m/s$^2$, $\rho=10^{3}$ kg/m$^3$, $T=\dfrac{1}{14}\times10^{-3}$ N/m)

\qn{2} The reading of the ideal ammeter connected in the given circuit is $0.39\times N$. The Value of $N$ is :- (Assume that the cells have negligible internal resistance)
\qfig{p1_phy_s2_q02_1}

\qn{3} The figure shows two identical parallel plate capacitors connected to a battery with switch S closed. The switch is now opened and the free space between the plates of the capacitors is filled with a dielectric of dielectric constant (or relative permittivity) 3. The ratio of the total electrostatic energy stored in both capacitors before and after the introduction of the dielectric is $3/N$. Then find the value of $N$.
\qfig{p1_phy_s2_q03_1}

\qn{4} A resistor of resistance $R$ and a capacitor of capacitance $C$ is connected in series with the voltage source. The EMF of the source changes periodically as shown in the figure. The maximum charge on the capacitor is $\dfrac{e}{e+N}\times CE$. The value of $N$ is -- ($\tau$ is time constant)
\qfig{p1_phy_s2_q04_1}

\qn{5} The space between the plates of a parallel plate capacitor is completely filled with a material of resistivity $2\times10^{11}\ \Omega$--m and dielectric constant 6. Capacity of the capacitor with the given dielectric medium between the plates is 20 nF. Find the leakage current (in $\mu$A) if a potential difference 2500 V is applied across the capacitor. (Give your answer in nearest integer.)

\qn{6} A tube of conical bore of length 10cm is just dipped inside the water. The diameters of upper and lower ends are 0.04 cm and 0.06 cm respectively. Find the height (in cm) to which the liquid rises in the tube. Given surface tension of water is 70 dyne cm$^{-1}$ and angle of contact $0^\circ$. (Take : $g=9.8$ m/s$^2$) (Give your answer in nearest integer.)

\qn{7} In the circuit shown in Figure,
\qfig{p1_phy_s2_q07_1}
current in the $3\ \Omega$ resistor is

\qn{8} A open organ pipe containing air resonates in fundamental mode due to a tuning fork. The measurement value of length $\ell$ (in cm) of the pipe are $\ell=94.0\pm0.1$, $r=5.0\pm0.05$. The velocity of the sound in air is accurately known. The maximum percentage error in the measurement of the frequency of that tuning fork by thus experiment, will be $\alpha$, then find the value of $25\alpha$.
\partsec{2}{CHEMISTRY}
\sect{SECTION-I(i)}
\qn{1} Consider following four reactions and given plot showing variation of equilibrium constant ($K_{eq}$) with temperature :

$\mathrm{A+B}\rightleftharpoons\mathrm{C}$\quad $\Delta H_1^\circ$

$\mathrm{D+E}\rightleftharpoons\mathrm{F}$\quad $\Delta H_2^\circ$

$\mathrm{G}\rightleftharpoons\mathrm{H}$\quad $\Delta H_3^\circ$

$\mathrm{I}\rightleftharpoons\mathrm{J}$\quad $\Delta H_4^\circ$
\qfig{p1_chem_s1_q01_1}
The correct graph showing the variation of \% yield of product, C, F, H and J with temperature is ($\Delta H^\circ$ is independent of temperature, $\ln K$ is the natural log of equilibrium constant)
\begin{qoptions}
\opt{A}{\qfigi{p1_chem_s1_q01_A}}
\opt{B}{\qfigi{p1_chem_s1_q01_B}}
\opt{C}{\qfigi{p1_chem_s1_q01_C}}
\opt{D}{\qfigi{p1_chem_s1_q01_D}}
\end{qoptions}

\qn{2} Consider the four different cells having equal cell potential is given conditions :

Ag $|$ Ag$^+$ (Ag$_2$CrO$_4$, saturated sol.) $\|$ Ag$^+$ (0.1M) $|$ Ag, $E_{cell}=x$V and $K_{sp,Ag_2CrO_4}=$ `p'

Ag $|$ Ag$^+$(AgCl saturated sol.) $\|$ Ag$^+$(0.1M) $|$ Ag, \quad $E_{cell}=x$V and $K_{sp,AgCl}=$ `q'

Ag $|$ Ag$^+$(Ag$_3$PO$_4$ saturated sol.) $\|$ Ag$^+$(0.1M) $|$ Ag, \quad $E_{cell}=x$V and $K_{sp,Ag_3PO_4}=$ `r'

Pt $|$ H$_2$(g), 1 bar $|$ H$^+$(1M) $\|$ (Ag$_2$SO$_4$ saturated sol.) Ag$^+$ $|$ Ag, \quad $E_{cell}=x$V and $K_{sp,Ag_2SO_4}=$ `s'

The correct option for solubility product for all the four cases is (by considering all cells are in equilibrium stage)
\begin{qoptions}
\opt{A}{$s>q>p>r$}
\opt{B}{$s<p<q<r$}
\opt{C}{$r>s>q>p$}
\opt{D}{$r>q>p>s$}
\end{qoptions}

\qn{3} 1 g of charcoal adsorbs 100 ml of 0.5 M CH$_3$COOH to form a monolayer, and thereby the molarity of CH$_3$COOH reduces to 0.49. Find the surface area of the charcoal adsorbed by each molecule of acetic acid. Surface area of charcoal $=3.01\times10^{2}$m$^2$/g.
\begin{qoptions}
\opt{A}{$5\times10^{-19}$ m$^2$}
\opt{B}{$9\times10^{-19}$ m$^2$}
\opt{C}{$20\times10^{-19}$ m$^2$}
\opt{D}{$25\times10^{-19}$ m$^2$}
\end{qoptions}

\qn{4} 10 g of a weak monocarboxylic acid in 100 g of benzene freezes at 3.35$^\circ$C. Assume that acid undergoes complete association in benzene. Find the degree of dissociation of 10 g of acid in 100 g of water whose freezing point is found to be $-1.75^\circ$C. Given: $K_f$ (H$_2$O) = 1.86 K-Kg/mole, $K_f$(benzene) = 5.12 K-Kg/mole and freezing point of benzene = 5.51$^\circ$C.
\begin{qoptions}
\opt{A}{0.05}
\opt{B}{0.115}
\opt{C}{0.621}
\opt{D}{0.825}
\end{qoptions}

\sect{SECTION-I(ii)}
\qn{5} The chromate ion, $\mathrm{CrO_4^{2-}}$ oxidizes the stannite ion, $\mathrm{HSnO_2^-}$ to the stannate ion $\mathrm{HSnO_3^-}$ and is itself reduced to chromite ion, $\mathrm{CrO_2^-}$ in basic medium as per the following reaction:
\[\mathrm{CrO_4^{2-}+HSnO_2^-+H_2O\to HSnO_3^-+CrO_2^-+OH^-}\]
If 400 ml of M/5 chromate ion react with 500 ml of M/4 stannite ion then:
\begin{qoptions}
\opt{A}{Chromate ion, $\mathrm{CrO_4^{2-}}$ is the limiting reagent}
\opt{B}{Stannite ion $\mathrm{HSnO_2^-}$ is the limiting reagent}
\opt{C}{At the end of reaction, the concentration of $\mathrm{CrO_2^-}$ ions $\approx0.09$M}
\opt{D}{At the end of reaction, the concentration of $\mathrm{HSnO_3^-}$ ions $\approx0.13$M}
\end{qoptions}

\qn{6} One mole of an ideal monoatomic gas initially at a temperature $T_1$, pressure $P_1$ and volume $V_1$ in a thermally insulated cylinder closed at one end and fitted with a movable piston at the other end is suddenly compressed placing a weight on the piston. The gas pressure is raised to $2P_1$. Which of the following is correct?
\begin{qoptions}
\opt{A}{Rise in temperature $=\Delta T=\dfrac{2}{5}T_1$}
\opt{B}{Change in volume $=\Delta V=\dfrac{-3}{10}\dfrac{RT_1}{P_1}$}
\opt{C}{Final temperature $=\dfrac{7}{5}T_1$}
\opt{D}{$P_1V_1^{\gamma}=(2P_1)V_2^{\gamma}$}
\end{qoptions}

\qn{7} A sample of H$_2$O$_2$ solution labelled as 34.05 volume has density of 1.128 gm/ml. Mark the correct option(s) representing concentration of same solution in other units. [Solution contains only H$_2$O and H$_2$O$_2$]
\begin{qoptions}
\opt{A}{Mole fraction of H$_2$O$_2$ in the solution $=0.5$}
\opt{B}{\%w/v $=10.2\%$}
\opt{C}{$M_{\mathrm{H_2O_2}}=6$M}
\opt{D}{$m_{\mathrm{H_2O_2}}=2.924$ m}
\end{qoptions}

\qn{8} For phase transition involving 3 moles of benzene --

C$_6$H$_6$($\ell$, 77$^\circ$C, 1 atm) $\to$ C$_6$H$_6$(g, 77$^\circ$C, 0.5 atm) ; $\Delta H_{vap}=1.4$ kJ/mole

If 77$^\circ$C is normal boiling point of benzene then correct value(s) of change in thermodynamic parameters for process is : [Use : $R=8$ J/K-mole , $\ln2=0.7$]
\begin{qoptions}
\opt{A}{$\Delta S=20.8$ J/K}
\opt{B}{$\Delta G=-5.88$ kJ}
\opt{C}{$\Delta U<0$}
\opt{D}{$\Delta U>0$}
\end{qoptions}

\qn{9} Which of following statement(s) is/are correct for Zeolites :
\begin{qoptions}
\opt{A}{They are hydrated alumino-silicates which can be used for shape selective catalysis}
\opt{B}{They have general formula $\mathrm{M_{x/n}[(AlO_2)_x(SiO_2)_y]}$\\$\mathrm{.mH_2O}$ where, n is the charge on the metal cation $\mathrm{M^{n+}}$}
\opt{C}{They have pore sizes between 260 pm to 740 pm}
\opt{D}{ZSM-5 which is used in petroleum industry as a catalyst, converts hydrocarbons into alcohols is an example of Zeolites.}
\end{qoptions}

\qn{10} Choose the correct statement(s) among the following --
\begin{qoptions}
\opt{A}{Internal energy of H$_2$O remains constant during conversion of liquid into its vapour at constant temperature}
\opt{B}{During fusion of ice into water, enthalpy change and internal energy change are almost same at constant temperature}
\opt{C}{Molar heat capacity of diatomic gases are temperature dependent}
\opt{D}{During boiling of H$_2$O at 1 atm, 100$^\circ$C, average kinetic energy of H$_2$O molecules increases}
\end{qoptions}

\sect{SECTION-III}
\qn{1} The variation of $\dfrac{1}{\lambda_m}$ vs $C\lambda_m$ for RCOOH at 298 K has been shown in the graph. Find p$K_a$ of RCOOH.
\qfig{p1_chem_s3_q01_1}

\qn{2} A 150 ml solution of I$_2$ is divided into two unequal parts. First part reacts with hypo solution in acidic medium. 15 ml of 0.4 M hypo was consumed. Second part was treated with 100 ml of 0.3M NaOH solution yielding NaI and NaIO$_3$. Residual base required 10 ml of 0.3M H$_2$SO$_4$ solution for complete neutralization. Calculate the value of ``10 times'' the initial concentration of I$_2$?

\qn{3} An aromatic compound contains 69.4\% carbon and 5.8\% hydrogen. A sample of 0.303 g of this compound was analyzed for nitrogen by kjeldahl's method. The ammonia evolved was absorbed in 50 ml of 0.05M H$_2$SO$_4$. The excess of the acid required 25 ml of 0.1 M NaOH for neutralization. Its molecular mass is 121. Determine the molecular formula and express the answer in the form of $\dfrac{w+x-y-z}{2}$ if the molecular formula of the compound is C$_w$H$_x$N$_y$O$_z$.

\qn{4} 0.00012\% MgSO$_4$ and 0.000111\% CaCl$_2$ is present in water. What is the measured hardness of water in ppm of CaCO$_3$. [Atomic mass : Mg = 24, S = 32, Ca = 40, Cl = 35.5]

\qn{5} When Aluminium pieces are added to aqueous solution of Tetrathionic acid (H$_2$S$_4$O$_6$) a pungent gas is produced. Find number of moles of Tetrathionic acid consumed per mole of Aluminium in this reaction ?

\qn{6} Calculate the quantity (in gram) of CO$_2$ in 550 mL of soda water when packed under 2.5 Bars CO$_2$ pressure at 298 K ($K_H$ for CO$_2=1.67\times10^{8}$ Pa at 298 K)

\qn{7} A 10 ml sample of hydroxylamine solution was diluted to one litre. 50 mL of this diluted solution was boiled with excess of Fe (III) solution and the resulting solution required 10 mL of 0.02M KMnO$_4$ solution for complete oxidation of Fe(II). Calculate the mass (in g) hydroxylamine in 91 mL of the original hydroxylamine solution.

Balanced reaction is, $2\mathrm{NH_2OH}+4\mathrm{Fe^{3+}}\to \mathrm{N_2O(g)}+\mathrm{H_2O}+4\mathrm{Fe^{2+}}+4\mathrm{H^+}$

\qn{8} In the following reactions:

$\mathrm{CaOCl_2}\xrightarrow{\Delta}a\,\mathrm{Ca(ClO_3)_2}+b\,\mathrm{CaCl_2}$

$\mathrm{NaCl+H_2SO_4\to Na_2SO_4}+c\,\mathrm{HCl}$

$\mathrm{Na_2S_2O_3}+d\,\mathrm{HCl}\to e\,\mathrm{NaCl+H_2O+S\downarrow+SO_2}$

What will be the value of $(a+b+c)-(d+e)$ in the balanced chemical reactions?
\partsec{3}{MATHEMATICS}
\sect{SECTION-I(i)}
\qn{1} The value(s) of the real parameter `a' for which the inequality $x^6-6x^5+12x^4+ax^3+12x^2-6x+1\ge0$ is satisfied for all real $x$
\begin{qoptions}
\opt{A}{$[-12,37]$}
\opt{B}{$[-12,38]$}
\opt{C}{$[-17,14]$}
\opt{D}{$[-17,35]$}
\end{qoptions}

\qn{2} Let $f(x)=\sin^3x+\lambda\sin^2x$ for $-\dfrac{\pi}{2}<x<\dfrac{\pi}{2}$ and $\lambda>0$. The interval in which $\lambda$ should lie in order that $f(x)$ has exactly one minimum and one maximum is
\begin{qoptions}
\opt{A}{$0<\lambda<\dfrac{3}{2}$}
\opt{B}{$0<\lambda<1$}
\opt{C}{$1<\lambda<2$}
\opt{D}{$1<\lambda<\dfrac{3}{2}$}
\end{qoptions}

\qn{3} The value of $\displaystyle\lim_{x\to2}\dfrac{\sqrt{x+7}-3\sqrt{2x-3}}{\sqrt[3]{x+6}-2\sqrt[3]{3x-5}}$ is equal to
\begin{qoptions}
\opt{A}{$\dfrac{34}{23}$}
\opt{B}{$-\dfrac{34}{23}$}
\opt{C}{$\dfrac{23}{34}$}
\opt{D}{$0$}
\end{qoptions}

\qn{4} The function defined as
\[f(x)=\lim_{n\to\infty}\begin{cases}\cos^{2n}x & \text{if } x<0\\[2pt] \sqrt[n]{1+x^n} & \text{if } 0\le x\le1\\[2pt] \dfrac{1}{1+x^n} & \text{if } x>1\end{cases}\]
which of the following does hold good.
\begin{qoptions}
\opt{A}{Continuous at $x=0$, but discontinuous at $x=1$}
\opt{B}{Continuous at $x=1$, but discontinuous at $x=0$}
\opt{C}{Continuous at both $x=1$ and $x=0$}
\opt{D}{Discontinuous at both $x=1$ and $x=0$}
\end{qoptions}

\sect{SECTION-I(ii)}
\qn{5} A function $f(x)$ is defined as
\[f(x)=\begin{cases}6x-5-x^2 & x\le3\\ 24x-32-4x^2 & x>3\end{cases}\]
Tangents are made on $f(x)$ at $(x,y)$ where $x\ge0$ \& $y\ge0$. Let tangents $T_1: y=m_1x+b_1$ and $T_2: y=m_2x+b_2$ have highest and lowest $y$ intercepts of all tangents made in first quadrant then
\begin{qoptions}
\opt{A}{$b_1+b_2=36$}
\opt{B}{$b_1+b_2=28$}
\opt{C}{$m_1+m_2=-4$}
\opt{D}{Area made by $y$ axis, $T_1=0$ and $T_2=0$ is 54 square units}
\end{qoptions}

\qn{6} Suppose that $f(x)$ is defined and positive for all real numbers $x$. If $f(x)$ increases for $x<0$ and decreases for all $x>0$, and $g(x)=\dfrac{1}{f\left(\frac1x\right)}$, then
\begin{qoptions}
\opt{A}{$g$ increases for all real $x$}
\opt{B}{$g$ increases for $x<0$}
\opt{C}{$g$ decreases for $x>0$}
\opt{D}{$g$ increases for $x>0$}
\end{qoptions}

\qn{7} Let $f(x)$ be double differentiable function such that $|f''(x)|\le5\ \forall\ x\in[0,4]$ and $f$ takes its largest value at an interior point of this interval. Then the value of $|f'(0)|+|f'(4)|$ can be
\begin{qoptions}
\opt{A}{18}
\opt{B}{19}
\opt{C}{20}
\opt{D}{21}
\end{qoptions}

\qn{8} If $F(x)=f(x)\,g(x)$ and $f'(x)\,g'(x)=c$, then
\begin{qoptions}
\opt{A}{$F'=c\left[\dfrac{f}{f'}+\dfrac{g}{g'}\right]$}
\opt{B}{$\dfrac{F''}{F}=\dfrac{f''}{f}+\dfrac{g''}{g}+\dfrac{2c}{fg}$}
\opt{C}{$\dfrac{F'''}{F}=\dfrac{f'''}{f}+\dfrac{g'''}{g}$}
\opt{D}{$\dfrac{F'''}{F''}=\dfrac{f'''}{f''}+\dfrac{g'''}{g''}$}
\end{qoptions}

\qn{9} Let $f:[0,1]\to\mathbb R$ be a double differentiable function with $f(0)=f(1)=0$, and $f''(x)+2f'(x)+f(x)\ge0\ \forall\ x\in(0,1)$, then
\begin{qoptions}
\opt{A}{$f(x)\le0\ \forall\ x\in[0,1]$}
\opt{B}{If $e^xf(x)$ assumes its minimum value in $[0,1]$ at $x=\dfrac{1}{\sqrt2}$, then $f(x)+f'(x)<0\ \forall\ x\in\left(0,\dfrac{1}{\sqrt2}\right)$}
\opt{C}{If $g(x)=e^xf(x)$, then number of real solution of the equation $g(g(g(x)))=0$ is 2}
\opt{D}{$f(x)>0\ \forall\ x\in[0,1]$}
\end{qoptions}

\qn{10} Which of the following are true for $\forall\ x\in(0,\infty)$
\begin{qoptions}
\opt{A}{$\ln(1+x)>x-\dfrac{x^2}{2}$}
\opt{B}{$\ln(1+x)>\dfrac{x}{x+1}$}
\opt{C}{$4\cos x+x>0$}
\opt{D}{$2\tan^{-1}x<x+1$}
\end{qoptions}

\sect{SECTION-III}
\qn{1} For the function\\$f(x)=\displaystyle\lim_{n\to\infty}\dfrac{\ln([1+9x^2])-9^nx^{2n}\cos(9x^2)}{1+9^nx^{2n}}$\\number of points of discontinuity in $(-3,3)$ is: (where $[.]$ represents greatest integer function)

\qn{2} Let $C$ be the curve $f(x)=(\ln x)^2+2\ln x$ and $A(a,f(a))$, $B(b,f(b))$ where $(a<b)$ are the points of tangency of two tangents drawn from origin to the curve $C$. The value of the product $ab$ is equal to

\qn{3} The smallest integral value of $P$ for which the function $f(x)=6Px-P\sin4x-5x-\sin3x$ is monotonically increasing and has no critical point on $R$ is:

\qn{4} Let $f(x)$ be a real valued function such that $f(x)=\dfrac{2x-1}{x-2}\ \forall\ x\in(2,\infty)$ and $g(x)=\dfrac{x^2+1}{x}+\dfrac{(f(x))^2+1}{f(x)}\ \forall\ x>2$, then minimum value of $g(x)$ is \underline{\hspace{2.5em}}.

\qn{5} Let $f:R\to R$ be a twice differentiable function such that $\displaystyle\lim_{x\to\infty}f(x)=a\ (a\ne0)$, then $\displaystyle\lim_{x\to\infty}x^2f''(x)$ is equal to

\qn{6} A continuous and differentiable function $y=f(x)$ is such that it's graph cuts line $y=mx+c$ at 11 distinct points. Then the minimum number of points at which $f''(x)=0$ is

\qn{7} The set of all points where $f(x)$ is increasing in $(a,b)\cup(c,\infty)$ then $[a+b+c]$ is (where $[.]$ denotes the greatest integer function). Given that $f(x)=2f\left(\dfrac{x^2}{2}\right)+f(6-x^2)\ \forall\ x\in R$ and $f''(x)>0\ \forall\ x\in R$

\qn{8} The equation $k+x^3e^{x+3}=0$ has 2 distinct real roots and $k$ is a prime natural number, then the number of possible value of $k$ is
\end{multicols}
\end{document}
